% ============================================================
\chapter{Mod\`eles GARCH et Volatilit\'e Conditionnelle}
\label{ch:garch}
% ============================================================

\section{Exemple motivant}
Les rendements financiers sont approximativement non corr\'el\'es, mais leur carr\'e $r_t^2$ montre une forte autocorr\'elation : les p\'eriodes de forte volatilit\'e sont group\'ees (\textbf{volatility clustering}\index{volatility clustering}). Les mod\`eles ARMA ne capturent pas ce ph\'enom\`ene.

\section{Faits stylis\'es des rendements financiers}\index{faits stylises@faits stylis\'es}

\begin{enumerate}
\item Rendements approximativement non corr\'el\'es : $\rho_{r}(h)\approx 0$.
\item Rendements au carr\'e corr\'el\'es : $\rho_{r^2}(h)>0$ (persistence de la volatilit\'e).
\item Queues lourdes (leptokurtose) : kurtosis $>3$.
\item Asym\'etrie de la volatilit\'e (effet de levier) : les baisses augmentent davantage la volatilit\'e que les hausses.
\end{enumerate}

\section{Mod\`ele ARCH}\index{ARCH}

\begin{definition}[ARCH$(q)$]
\[
r_t = \sigma_t\,z_t, \quad z_t\iid(0,1), \quad \sigma_t^2 = \omega + \sum_{i=1}^q\alpha_i r_{t-i}^2,
\]
o\`u $\omega>0$, $\alpha_i\geq 0$, et $\sum\alpha_i<1$ pour la stationnarit\'e.
\end{definition}

\begin{theoreme}[Propri\'et\'es d'un ARCH$(q)$]
\begin{enumerate}
\item $\E[r_t]=0$, $\Var(r_t)=\omega/(1-\sum\alpha_i)$.
\item $(r_t)$ est un bruit blanc (non corr\'el\'e).
\item $(r_t^2)$ suit un AR$(q)$ : $r_t^2 = \omega+\sum\alpha_i r_{t-i}^2+v_t$ o\`u $v_t=r_t^2-\sigma_t^2$ est un bruit blanc.
\item Kurtosis $>3$ (queues lourdes m\^eme si $z_t$ est gaussien).
\end{enumerate}
\end{theoreme}

\section{Mod\`ele GARCH}\index{GARCH}

\begin{definition}[GARCH$(p,q)$]
\[
r_t = \sigma_t z_t, \quad \sigma_t^2 = \omega + \sum_{i=1}^q\alpha_i r_{t-i}^2 + \sum_{j=1}^p\beta_j\sigma_{t-j}^2.
\]
Conditions de stationnarit\'e : $\omega>0$, $\alpha_i,\beta_j\geq 0$, $\sum_{i=1}^q\alpha_i+\sum_{j=1}^p\beta_j<1$.
\end{definition}

\begin{theoreme}[Variance inconditionnelle GARCH$(1,1)$]
\[
\Var(r_t) = \frac{\omega}{1-\alpha_1-\beta_1}.
\]
La \textbf{persistence} de la volatilit\'e est mesur\'ee par $\alpha_1+\beta_1$ (proche de 1 en finance).
\end{theoreme}

\begin{intuition}
Le GARCH$(1,1)$ est le mod\`ele de volatilit\'e le plus utilis\'e en pratique. La pr\'evision de la variance conditionnelle un pas en avant est :
\[
\hat{\sigma}_{t+1}^2 = \omega + \alpha_1 r_t^2 + \beta_1\sigma_t^2.
\]
C'est une moyenne pond\'er\'ee exponentielle des carr\'es pass\'es.
\end{intuition}

\section{IGARCH}\index{IGARCH}

\begin{definition}
Si $\alpha_1+\beta_1=1$ (int\'egr\'e GARCH), les chocs de volatilit\'e persistent ind\'efiniment. Le processus n'est pas stationnaire au second ordre mais reste \textbf{strictement stationnaire}.
\end{definition}

\section{Extensions}

\subsection{EGARCH}\index{EGARCH}
\begin{definition}
Le mod\`ele EGARCH de Nelson :
\[
\ln\sigma_t^2 = \omega + \alpha(|z_{t-1}|-\E[|z_{t-1}|]) + \gamma z_{t-1} + \beta\ln\sigma_{t-1}^2.
\]
Le terme $\gamma z_{t-1}$ capture l'\textbf{effet de levier}. Pas de contrainte de positivit\'e sur les param\`etres.
\end{definition}

\subsection{GJR-GARCH}\index{GJR-GARCH}
\begin{definition}
\[
\sigma_t^2 = \omega + (\alpha+\gamma\mathbf{1}_{r_{t-1}<0})r_{t-1}^2 + \beta\sigma_{t-1}^2.
\]
$\gamma>0$ implique que les rendements n\'egatifs augmentent davantage la volatilit\'e.
\end{definition}

\section{Estimation}\index{estimation!GARCH}

\begin{definition}[Quasi-maximum de vraisemblance]
La log-vraisemblance conditionnelle gaussienne :
\[
\ell(\bm\theta) = -\frac{1}{2}\sum_{t=1}^n\left(\ln\sigma_t^2 + \frac{r_t^2}{\sigma_t^2}\right).
\]
M\^eme si $z_t$ n'est pas gaussien, le QMLE est consistant sous certaines conditions.
\end{definition}

\section{Tests d'effet ARCH}\index{test!ARCH}

\begin{definition}[Test ARCH-LM d'Engle]
R\'egresser $r_t^2$ sur $r_{t-1}^2,\ldots,r_{t-q}^2$ et tester la nullit\'e jointe par $nR^2\sim\chi^2(q)$.
\end{definition}

\section{Impl\'ementation}

\begin{codeblock}[title=\textbf{Python}]
\begin{minted}{python}
import numpy as np
import pandas as pd
import matplotlib.pyplot as plt
from arch import arch_model

# Simuler un GARCH(1,1)
np.random.seed(42)
n = 1000
omega, alpha, beta = 0.05, 0.10, 0.85
sigma2 = np.zeros(n)
r = np.zeros(n)
sigma2[0] = omega / (1 - alpha - beta)
for t in range(1, n):
    sigma2[t] = omega + alpha * r[t-1]**2 + beta * sigma2[t-1]
    r[t] = np.sqrt(sigma2[t]) * np.random.normal()

fig, axes = plt.subplots(3, 1, figsize=(12, 8))
axes[0].plot(r, lw=0.5)
axes[0].set_title('Rendements simulés (GARCH(1,1))')
axes[1].plot(r**2, lw=0.5, color='orange')
axes[1].set_title('$r_t^2$ (volatilité réalisée)')
axes[2].plot(np.sqrt(sigma2), lw=0.5, color='red')
axes[2].set_title('$\\sigma_t$ (volatilité conditionnelle)')
plt.tight_layout()
plt.savefig('ch06_garch.pdf')
plt.show()

# Ajuster un GARCH(1,1)
am = arch_model(r, vol='Garch', p=1, q=1, dist='normal')
res = am.fit(disp='off')
print(res.summary())

# Prévision de volatilité
forecasts = res.forecast(horizon=10)
print("Variance prévisionnelle:")
print(forecasts.variance.iloc[-1])
\end{minted}
\end{codeblock}

\section{Exercices}

\begin{exercice}[$\star$]
Montrer qu'un ARCH(1) avec $r_t=\sigma_t z_t$, $\sigma_t^2=\omega+\alpha r_{t-1}^2$ implique $\E[r_t^4]=3\omega^2(1+\alpha)/(1-3\alpha^2)$ (si $3\alpha^2<1$). En d\'eduire la kurtosis.
\end{exercice}

\begin{exercice}[$\star\star$]
Montrer que le GARCH$(1,1)$ peut s'\'ecrire comme un ARCH$(\infty)$.
\end{exercice}

\begin{exercice}[$\star\star$]
D\'emontrer que le processus $r_t^2$ d'un GARCH$(1,1)$ suit un ARMA$(1,1)$.
\end{exercice}

\begin{exercice}[$\star\star\star$ -- Projet]
T\'el\'echarger les rendements journaliers du S\&P~500 sur 10 ans. Estimer un GARCH$(1,1)$, un EGARCH$(1,1)$ et un GJR-GARCH$(1,1)$. Comparer par AIC et analyser l'effet de levier.
\end{exercice}

\begin{formules}
\begin{align*}
\text{GARCH}(1,1) &: \sigma_t^2 = \omega+\alpha r_{t-1}^2+\beta\sigma_{t-1}^2, \quad \alpha+\beta<1 \\
\Var(r_t) &= \omega/(1-\alpha-\beta) \\
\text{EGARCH} &: \ln\sigma_t^2 = \omega+\alpha(|z_{t-1}|-\E|z_{t-1}|)+\gamma z_{t-1}+\beta\ln\sigma_{t-1}^2
\end{align*}
\end{formules}
